\documentclass[11pt]{article}

\usepackage[a4paper,margin=1in]{geometry}
\usepackage{amsmath,amssymb,amsthm,mathtools}
\usepackage[T1]{fontenc}
\usepackage{lmodern}
\usepackage{microtype}
\usepackage{hyperref}

\newtheorem{theorem}{Theorem}
\newtheorem{lemma}[theorem]{Lemma}
\newtheorem{remark}[theorem]{Remark}

\DeclareMathOperator{\tridiag}{tridiag}
\DeclareMathOperator{\diag}{diag}
\DeclareMathOperator{\Res}{Res}
\newcommand{\A}{\mathbf A}
\newcommand{\I}{\mathbf I}
\newcommand{\R}{\mathbb R}
\newcommand{\C}{\mathbb C}
\newcommand{\eps}{\varepsilon}

\begin{document}

\begin{theorem}[Exact algebraic formula for the connectedness threshold]
Let \(0<a<1\) and \(\eps>0\). For \(n\ge2\), put
\[
\A_n(a)=\tridiag(a,0,1).
\]
Define the increasing ordering of the eigenvalues of \(\A_n(a)\) by
\[
\lambda_{n,k}(a)
=
2\sqrt a\cos\frac{(n+1-k)\pi}{n+1},
\qquad k=1,\dots,n.
\]
Thus
\[
\lambda_{n,1}(a)<\lambda_{n,2}(a)<\cdots<\lambda_{n,n}(a).
\]

For \(x\in\R\), define
\[
M_n(x,a):=(x\I_n-\A_n(a))^T(x\I_n-\A_n(a)).
\]
For \(u\in\R\), set
\[
Q_n(x,u;a):=M_n(x,a)-u\I_n.
\]
Equivalently, \(Q_n(x,u;a)=[q_{ij}]_{i,j=1}^n\) is the real symmetric pentadiagonal matrix determined by
\[
q_{jj}
=
x^2-u+
\begin{cases}
a^2, & j=1,\\
1+a^2, & 2\le j\le n-1,\\
1, & j=n,
\end{cases}
\]
\[
q_{j,j+1}=q_{j+1,j}=-(1+a)x,
\qquad 1\le j\le n-1,
\]
\[
q_{j,j+2}=q_{j+2,j}=a,
\qquad 1\le j\le n-2,
\]
and all other entries are zero.

For \(1\le k\le n-1\), define
\[
\Gamma_{n,k}(a)
:=
\max
\left\{
u\ge0:
\begin{array}{l}
\text{there exists }x\in[\lambda_{n,k}(a),\lambda_{n,k+1}(a)]\\[2mm]
\text{such that }Q_n(x,u;a)\succeq0
\text{ and }\det Q_n(x,u;a)=0
\end{array}
\right\}.
\]
Here \(Q\succeq0\) means that \(Q\) is positive semidefinite. Finally set
\[
\Gamma_n(a):=\max_{1\le k\le n-1}\Gamma_{n,k}(a).
\]

Then the connectedness threshold is exactly
\[
\boxed{
N_c(a,\eps)
=
\min\left\{
n\ge2:\ \eps^2>\Gamma_n(a)
\right\}.
}
\]
Equivalently,
\[
\boxed{
N_c(a,\eps)
=
\min\left\{
n\ge2:
\eps^2>
\max_{1\le k\le n-1}
\max_{\lambda_{n,k}(a)\le x\le\lambda_{n,k+1}(a)}
s_{\min}\!\left(x\I_n-\A_n(a)\right)^2
\right\}.
}
\]
This is an exact finite algebraic formula. Indeed, for fixed \(n\), the condition
\[
Q_n(x,u;a)\succeq0
\]
is equivalent to the non-negativity of all principal minors of the explicitly given matrix \(Q_n(x,u;a)\). Hence each \(\Gamma_{n,k}(a)\) is the maximum of \(u\) over a compact semialgebraic set defined only by polynomial equations and inequalities in \(a,x,u\).
\end{theorem}

\begin{proof}
We prove the formula carefully.

First, the spectrum of \(\A_n(a)\) is explicit. Let
\[
D_n=\diag(1,a^{1/2},a,\dots,a^{(n-1)/2}).
\]
Then
\[
D_n^{-1}\A_n(a)D_n
=
\sqrt a\,\tridiag(1,0,1).
\]
The symmetric tridiagonal matrix \(\tridiag(1,0,1)\) has eigenvalues
\[
2\cos\frac{j\pi}{n+1},
\qquad j=1,\dots,n.
\]
Therefore
\[
\sigma(\A_n(a))
=
\left\{
2\sqrt a\cos\frac{j\pi}{n+1}:j=1,\dots,n
\right\},
\]
which is exactly the ordered list \(\lambda_{n,1}(a)<\cdots<\lambda_{n,n}(a)\) above.

Now define
\[
g_n(x,a):=
s_{\min}\!\left(x\I_n-\A_n(a)\right)^2
=
\lambda_{\min}\!\left(M_n(x,a)\right).
\]
For fixed \(x\), the set
\[
\left\{
u\ge0:Q_n(x,u;a)\succeq0,\ \det Q_n(x,u;a)=0
\right\}
\]
contains exactly one number, namely \(u=g_n(x,a)\). Indeed, if
\[
\mu_1\le\mu_2\le\cdots\le\mu_n
\]
are the eigenvalues of \(M_n(x,a)\), then
\[
Q_n(x,u;a)=M_n(x,a)-u\I_n\succeq0
\]
is equivalent to \(u\le\mu_1\), while
\[
\det Q_n(x,u;a)=0
\]
is equivalent to \(u=\mu_j\) for at least one \(j\). Both conditions together force
\[
u=\mu_1=\lambda_{\min}(M_n(x,a))=g_n(x,a).
\]
Conversely, \(u=g_n(x,a)\) plainly satisfies both conditions. Hence
\[
\Gamma_{n,k}(a)
=
\max_{\lambda_{n,k}(a)\le x\le\lambda_{n,k+1}(a)}
g_n(x,a).
\]
Consequently,
\[
\Gamma_n(a)
=
\max_{1\le k\le n-1}
\max_{\lambda_{n,k}(a)\le x\le\lambda_{n,k+1}(a)}
s_{\min}\!\left(x\I_n-\A_n(a)\right)^2.
\]

We next recall the real-axis reduction for this Toeplitz family. Since
\[
\A_n(a)^TJ_n=J_n\A_n(a),
\]
where \(J_n\) is the reversal identity matrix, the matrix
\[
J_n(x\I_n-\A_n(a))
\]
is real symmetric for every \(x\in\R\). The singular-value equation for
\[
(x+iy)\I_n-\A_n(a)
\]
then reduces, after pairing the \(j\)-th and \((n+1-j)\)-th coordinates, to a positive second-order Jacobi recurrence. In that recurrence the parameter \(y\) enters through the positive term \(y^2\). Hence
\[
s_{\min}\!\left((x+iy)\I_n-\A_n(a)\right)
\ge
s_{\min}\!\left(x\I_n-\A_n(a)\right)
\qquad (x,y\in\R).
\]
Therefore every vertical section of the pseudospectrum
\[
\sigma_\eps(\A_n(a))
=
\left\{
z\in\C:
s_{\min}(z\I_n-\A_n(a))<\eps
\right\}
\]
is either empty or an open interval symmetric about the real axis. In particular, \(\sigma_\eps(\A_n(a))\) deformation retracts onto its real trace
\[
E_{n,\eps}(a)
:=
\sigma_\eps(\A_n(a))\cap\R
=
\left\{
x\in\R:g_n(x,a)<\eps^2
\right\}.
\]
Thus
\[
\sigma_\eps(\A_n(a))\ \text{is connected}
\quad\Longleftrightarrow\quad
E_{n,\eps}(a)\ \text{is connected}.
\]

The function \(g_n(\cdot,a)\) is continuous, nonnegative, and satisfies
\[
g_n(x,a)=0
\quad\Longleftrightarrow\quad
x\in\sigma(\A_n(a)).
\]
Hence \(E_{n,\eps}(a)\) contains a neighborhood of every eigenvalue
\[
\lambda_{n,1}(a),\dots,\lambda_{n,n}(a).
\]
The only possible disconnections of \(E_{n,\eps}(a)\) occur inside the real gaps
\[
[\lambda_{n,k}(a),\lambda_{n,k+1}(a)],
\qquad k=1,\dots,n-1.
\]
The whole \(k\)-th gap is contained in \(E_{n,\eps}(a)\) if and only if
\[
g_n(x,a)<\eps^2
\qquad
\text{for every }
x\in[\lambda_{n,k}(a),\lambda_{n,k+1}(a)].
\]
Because the maximum of \(g_n\) on this compact interval is \(\Gamma_{n,k}(a)\), this is equivalent to
\[
\Gamma_{n,k}(a)<\eps^2.
\]
Therefore
\[
E_{n,\eps}(a)\ \text{is connected}
\quad\Longleftrightarrow\quad
\Gamma_{n,k}(a)<\eps^2
\ \text{for every }k=1,\dots,n-1,
\]
or equivalently,
\[
E_{n,\eps}(a)\ \text{is connected}
\quad\Longleftrightarrow\quad
\Gamma_n(a)<\eps^2.
\]
The strict inequality is essential, because the pseudospectrum is defined by
\[
s_{\min}(z\I_n-\A_n(a))<\eps.
\]
If \(\Gamma_n(a)=\eps^2\), then at least one real gap contains a point at which
\[
g_n(x,a)=\eps^2,
\]
and that point is not included in the real trace. Hence the real trace is still disconnected.

Combining the preceding equivalences gives
\[
\sigma_\eps(\A_n(a))\ \text{is connected}
\quad\Longleftrightarrow\quad
\eps^2>\Gamma_n(a).
\]
Since by definition
\[
N_c(a,\eps)=N_f(a,\eps)
=
\min\left\{
n\ge2:
\sigma_\eps(\A_n(a))\ \text{is connected}
\right\},
\]
we obtain
\[
N_c(a,\eps)
=
\min\left\{
n\ge2:\eps^2>\Gamma_n(a)
\right\}.
\]

It remains only to note that the minimum is finite. For
\[
x=2\sqrt a\cos\theta,\qquad 0\le\theta\le\pi,
\]
take the vector
\[
q^{(n)}(\theta)
=
\left(
\sin\theta,\,
a^{1/2}\sin2\theta,\,
a\sin3\theta,\,
\dots,\,
a^{(n-1)/2}\sin n\theta
\right)^T.
\]
Using
\[
2\cos\theta\sin(j\theta)
=
\sin((j+1)\theta)+\sin((j-1)\theta),
\]
one obtains
\[
(x\I_n-\A_n(a))q^{(n)}(\theta)
=
a^{n/2}\sin((n+1)\theta)e_n.
\]
Therefore, for \(x\in[-2\sqrt a,2\sqrt a]\),
\[
s_{\min}(x\I_n-\A_n(a))
\le
\frac{
a^{n/2}|\sin((n+1)\theta)|
}{
\left(\sum_{j=1}^n a^{j-1}\sin^2(j\theta)\right)^{1/2}
}
\le
(n+1)a^{n/2}.
\]
All spectral gaps lie in \([-2\sqrt a,2\sqrt a]\). Hence
\[
0\le\Gamma_n(a)\le (n+1)^2a^n\longrightarrow0
\qquad(n\to\infty),
\]
because \(0<a<1\). Thus for every fixed \(\eps>0\), there exists \(n\) such that
\[
\eps^2>\Gamma_n(a),
\]
so the displayed minimum is finite. The formula is proved.
\end{proof}

\begin{remark}
For \(n=2\), the formula gives
\[
\Gamma_2(a)=a^2,
\]
and therefore
\[
\sigma_\eps(\A_2(a))\ \text{is connected}
\quad\Longleftrightarrow\quad
\eps>a.
\]
For \(n=3\), writing
\[
C(a):=27a^4+36a^3+2a^2+36a+27,
\]
one obtains
\[
\Gamma_3(a)
=
\frac{
C(a)-\sqrt{C(a)^2-4096a^3(1+a)^2}
}{
64(1+a)^2
}.
\]
For higher \(n\), the same determinant formula with \(Q_n(x,u;a)\) gives the exact algebraic value of \(\Gamma_n(a)\).
\end{remark}

\end{document}